\section{Lowest Common Ancestor in a BST}
\label{sec:bst-lca}

\problemstatement{Given the root of a BST and two nodes $p$ and $q$
(both guaranteed to exist in the tree), find their \textbf{lowest common
ancestor} (LCA): the deepest node that has both $p$ and $q$ as
descendants (a node is considered a descendant of itself).}

\begin{patternbox}
LCA in a general binary tree requires examining both subtrees at every
node and combining results bottom-up, costing $O(n)$. The keyword
\emph{``BST''} attached to an LCA problem is a strong signal that the
BST invariant alone determines the split point directly, top-down, in
$O(h)$ -- no need to inspect both subtrees blindly the way a general
tree's LCA algorithm would.
\end{patternbox}

\subsection*{Understanding the problem carefully}

In a BST, the LCA of $p$ and $q$ is exactly the node where the search
paths to $p$ and $q$ \textbf{diverge}: as long as both $p.\textit{val}$
and $q.\textit{val}$ are on the same side (both less than the current
node, or both greater), the current node cannot be the LCA -- both
targets are reachable only by continuing in that single, same direction,
so the LCA must lie further down on that side. The moment $p$ and $q$
fall on \emph{different} sides (or one of them \emph{is} the current
node), the current node is exactly the point where their paths separate
-- the LCA.

\begin{ideabox}[Follow Both Targets Together Until They Diverge]
Starting at the root: if both $p.\textit{val}$ and $q.\textit{val}$ are
less than the current node's value, move left. If both are greater, move
right. Otherwise (one is $\le$ and the other is $\ge$ the current node's
value, including the case where the current node \emph{is} $p$ or $q$),
the current node is the LCA.
\end{ideabox}

\subsection*{Why the divergence point is provably the LCA}

If both targets are smaller than the current node, the current node
cannot be an ancestor that is also the \emph{lowest} one -- both $p$ and
$q$ are confirmed to live in the left subtree, so a deeper common
ancestor must exist there, and continuing left can only get closer to it
(never skip past it, since the BST invariant guarantees the entire left
subtree is exactly where both targets are). The symmetric argument holds
for both being larger. The first node where this is no longer true --
where the two targets are no longer both on the same side -- is, by
definition, the deepest point both search paths still share, since any
node below it would have to continue toward only one of $p$ or $q$, not
both.

\subsection*{Algorithm}

\begin{enumerate}
  \item Start at the root.
  \item While true:
  \begin{enumerate}
    \item If $p.\textit{val} < \textit{current}.\textit{val}$ and
          $q.\textit{val} < \textit{current}.\textit{val}$: move to the
          left child.
    \item Else if $p.\textit{val} > \textit{current}.\textit{val}$ and
          $q.\textit{val} > \textit{current}.\textit{val}$: move to the
          right child.
    \item Else: return \textit{current} (this is the LCA).
  \end{enumerate}
\end{enumerate}

\subsection*{Dry run}

\dryrun{Take the BST with root $6$, left child $2$ (children $0,4$),
right child $8$ (children $7,9$). Find LCA of $p=2$, $q=4$.}

\begin{itemize}
  \item At $6$: is $2<6$ and $4<6$? Yes, both smaller. Move left to
        $2$.
  \item At $2$: is $2<2$? No ($p$ equals the current node). So not both
        strictly smaller, and not both strictly larger. Return $2$.
\end{itemize}

LCA is $2$ (correct: $2$ is an ancestor of $4$, and is itself $p$).

\subsection*{C++ implementation}

\begin{lstlisting}
#include <bits/stdc++.h>
using namespace std;

struct TreeNode {
    int val;
    TreeNode* left;
    TreeNode* right;
    TreeNode(int v) : val(v), left(nullptr), right(nullptr) {}
};

TreeNode* lowestCommonAncestor(TreeNode* root, TreeNode* p, TreeNode* q) {
    TreeNode* current = root;
    while (true) {
        if (p->val < current->val && q->val < current->val) {
            current = current->left;
        } else if (p->val > current->val && q->val > current->val) {
            current = current->right;
        } else {
            return current;
        }
    }
}

int main() {
    TreeNode* root = new TreeNode(6);
    root->left = new TreeNode(2);
    root->right = new TreeNode(8);
    root->left->left = new TreeNode(0);
    root->left->right = new TreeNode(4);
    root->right->left = new TreeNode(7);
    root->right->right = new TreeNode(9);

    TreeNode* lca = lowestCommonAncestor(root, root->left, root->left->right);
    cout << "LCA: " << lca->val << endl;
    return 0;
}
\end{lstlisting}

\begin{complexitybox}
\textbf{Time Complexity.} $O(h)$ -- a single root-to-LCA walk, never
touching any subtree that doesn't lie on the path to the answer.
\textbf{Space Complexity.} $O(1)$ for the iterative version shown.
\end{complexitybox}

\begin{takeawaybox}
LCA in a general binary tree needs $O(n)$ because there is no way to know
which subtree contains which target without searching both. The BST
invariant removes this ambiguity entirely: comparing both targets'
values to the current node tells you, with certainty, which single
direction to go (or that you have already arrived), collapsing an
$O(n)$ two-subtree search into an $O(h)$ single-path walk.
\end{takeawaybox}
